% GENERATED FILE. Edit canonical lessons, then run npm run generate:books.
% canonical-source-hash: fnv1a64:bead3385d5cdebc7
% canonical-lessons: UR-C396-shahrah, UR-C396-mahsul, UR-C396-naka, UR-C396-qafila, UR-C396-hankna

\chapter{On the Highway}
\label{ch:ur-on-the-highway}
\hlchaptermodality{396}

\begin{chapteropening}
\textbf{By the end of this chapter:} \emph{I can talk about the highway: tolls, checkpoints and convoys of lorries.}

Complete the last lesson of chapter 396: I can talk about the highway: tolls, checkpoints and convoys of lorries.
\end{chapteropening}

\section[\texorpdfstring{\ur{شاہراہ} (shāhrāh)}{شاہراہ (shāhrāh)}]{\ur{شاہراہ} (shāhrāh) — a highway; a main road}

\label{lesson:UR-C396-shahrah}



\begin{quote}
Before the new one: say the Urdu for a helmet, then the Urdu for to get dented.
\end{quote}

\subsection*{You'll want to know: \ur{شاہراہ}}
\textbf{\ur{شاہراہ}} — \emph{shāhrāh} — \textquotedblleft{}a highway; a main road\textquotedblright{}.

\begin{quote}
\ur{شاہراہ} \ur{پر} \ur{ٹرک} \ur{تیز} \ur{چلتے} \ur{ہیں۔}
\end{quote}

\begin{quote}
\emph{shāhrāh par ṭrak tez chalte haiṅ.}
\end{quote}

\begin{quote}
Lorries drive fast on the highway.
\end{quote}

Feminine.

\begin{cousinweb}
\textbf{\ur{شاہراہ}} (\emph{shāhrāh}) is built from Persian \emph{shāh}, \textquotedblleft{}a king\textquotedblright{}, and \emph{rāh}, \textquotedblleft{}a road\textquotedblright{} (\textbf{\ur{راہ}}): it is usually traced to Persian \emph{shāhrāh}, \textquotedblleft{}the king's road\textquotedblright{}.
\end{cousinweb}

\begin{grammarlens}[title={What you've built}]
One more word for getting around.
\end{grammarlens}

\subsection*{Your turn}
\begin{itemize}
\raggedright
  \item \emph{Say it:} \emph{shāhrāh}
  \item \emph{Say it:} \emph{shāhrāh}, once more
  \item \emph{Say it:} \emph{shāhrāh par ṭrak tez chalte haiṅ.} — Lorries drive fast on the highway.
  \item \emph{Recall:} say the Urdu for a helmet, then the Urdu for to get dented, then say \emph{shāhrāh} again
\end{itemize}

\begin{tcolorbox}[breakable,colback=teal!4,colframe=teal!35!black,title={Before you move on}]
What does \ur{شاہراہ} mean? (\textquotedblleft{}A highway; a main road\textquotedblright{}.) Say it once more.
\end{tcolorbox}
\section[\texorpdfstring{\ur{محصول} (mahsūl)}{محصول (mahsūl)}]{\ur{محصول} (mahsūl) — a toll; a duty (on goods)}

\label{lesson:UR-C396-mahsul}



\begin{quote}
Before the new one: say the Urdu for to get dented, then the Urdu for a highway.
\end{quote}

\subsection*{You'll want to know: \ur{محصول}}
\textbf{\ur{محصول}} — \emph{mahsūl} — \textquotedblleft{}a toll; a duty (on goods)\textquotedblright{}.

\begin{quote}
\ur{شاہراہ} \ur{پر} \ur{ہر} \ur{گاڑی} \ur{محصول} \ur{دیتی} \ur{ہے۔}
\end{quote}

\begin{quote}
\emph{shāhrāh par har gāṛī mahsūl detī hai.}
\end{quote}

\begin{quote}
Every vehicle pays a toll on the highway.
\end{quote}

Masculine.

\begin{cousinweb}
\textbf{\ur{محصول}} (\emph{mahsūl}) is usually traced, through Persian, to Arabic \emph{maḥṣūl}, \textquotedblleft{}what is gathered, revenue\textquotedblright{}; in Urdu it is a toll, or a duty paid on goods.
\end{cousinweb}

\begin{grammarlens}[title={What you've built}]
One more word for getting around.
\end{grammarlens}

\subsection*{Your turn}
\begin{itemize}
\raggedright
  \item \emph{Say it:} \emph{mahsūl}
  \item \emph{Say it:} \emph{mahsūl}, once more
  \item \emph{Say it:} \emph{shāhrāh par har gāṛī mahsūl detī hai.} — Every vehicle pays a toll on the highway.
  \item \emph{Recall:} say the Urdu for to get dented, then the Urdu for a highway, then say \emph{mahsūl} again
\end{itemize}

\begin{tcolorbox}[breakable,colback=teal!4,colframe=teal!35!black,title={Before you move on}]
What does \ur{محصول} mean? (\textquotedblleft{}A toll; a duty (on goods)\textquotedblright{}.) Say it once more.
\end{tcolorbox}
\section[\texorpdfstring{\ur{ناکہ} (nākā)}{ناکہ (nākā)}]{\ur{ناکہ} (nākā) — a checkpoint (on a road)}

\label{lesson:UR-C396-naka}



\begin{quote}
Before the new one: say the Urdu for a highway, then the Urdu for a toll.
\end{quote}

\subsection*{You'll want to know: \ur{ناکہ}}
\textbf{\ur{ناکہ}} — \emph{nākā} — \textquotedblleft{}a checkpoint (on a road)\textquotedblright{}.

\begin{quote}
\ur{پولیس} \ur{کا} \ur{ناکہ} \ur{پل} \ur{کے} \ur{پاس} \ur{ہے۔}
\end{quote}

\begin{quote}
\emph{polīs kā nākā pul ke pās hai.}
\end{quote}

\begin{quote}
The police checkpoint is near the bridge.
\end{quote}

Masculine.

\begin{cousinweb}
\textbf{\ur{ناکہ}} (\emph{nākā}) is inherited: CDIAL files it under \emph{nākka}, \textquotedblleft{}like a nose, projecting\textquotedblright{}, a form it reconstructs (CDIAL 7037); the sense \textquotedblleft{}a checkpoint\textquotedblright{} is probably a later one, from a post that juts into the road.
\end{cousinweb}

\begin{grammarlens}[title={What you've built}]
One more word for getting around.
\end{grammarlens}

\subsection*{Your turn}
\begin{itemize}
\raggedright
  \item \emph{Say it:} \emph{nākā}
  \item \emph{Say it:} \emph{nākā}, once more
  \item \emph{Say it:} \emph{polīs kā nākā pul ke pās hai.} — The police checkpoint is near the bridge.
  \item \emph{Recall:} say the Urdu for a highway, then the Urdu for a toll, then say \emph{nākā} again
\end{itemize}

\begin{tcolorbox}[breakable,colback=teal!4,colframe=teal!35!black,title={Before you move on}]
What does \ur{ناکہ} mean? (\textquotedblleft{}A checkpoint (on a road)\textquotedblright{}.) Say it once more.
\end{tcolorbox}
\section[\texorpdfstring{\ur{قافلہ} (qāfilā)}{قافلہ (qāfilā)}]{\ur{قافلہ} (qāfilā) — a convoy; a caravan}

\label{lesson:UR-C396-qafila}



\begin{quote}
Before the new one: say the Urdu for a toll, then the Urdu for a checkpoint (on a road).
\end{quote}

\subsection*{You'll want to know: \ur{قافلہ}}
\textbf{\ur{قافلہ}} — \emph{qāfilā} — \textquotedblleft{}a convoy; a caravan\textquotedblright{}.

\begin{quote}
\ur{ناکے} \ur{پر} \ur{ٹرکوں} \ur{کا} \ur{ایک} \ur{قافلہ} \ur{رکا۔}
\end{quote}

\begin{quote}
\emph{nāke par ṭrakoṅ kā ek qāfilā rukā.}
\end{quote}

\begin{quote}
A convoy of lorries stopped at the checkpoint.
\end{quote}

Masculine.

\begin{cousinweb}
\textbf{\ur{قافلہ}} (\emph{qāfilā}) is usually traced, through Persian, to Arabic \emph{qāfila}, \textquotedblleft{}a caravan\textquotedblright{}.
\end{cousinweb}

\begin{grammarlens}[title={What you've built}]
One more word for getting around.
\end{grammarlens}

\subsection*{Your turn}
\begin{itemize}
\raggedright
  \item \emph{Say it:} \emph{qāfilā}
  \item \emph{Say it:} \emph{qāfilā}, once more
  \item \emph{Say it:} \emph{nāke par ṭrakoṅ kā ek qāfilā rukā.} — A convoy of lorries stopped at the checkpoint.
  \item \emph{Recall:} say the Urdu for a toll, then the Urdu for a checkpoint (on a road), then say \emph{qāfilā} again
\end{itemize}

\begin{tcolorbox}[breakable,colback=teal!4,colframe=teal!35!black,title={Before you move on}]
What does \ur{قافلہ} mean? (\textquotedblleft{}A convoy; a caravan\textquotedblright{}.) Say it once more.
\end{tcolorbox}
\section[\texorpdfstring{\ur{ہانکنا} (hāṅknā)}{ہانکنا (hāṅknā)}]{\ur{ہانکنا} (hāṅknā) — to drive (animals or a cart)}

\label{lesson:UR-C396-hankna}



\begin{quote}
Before the new one: say the Urdu for a checkpoint (on a road), then the Urdu for a convoy.
\end{quote}

\subsection*{You'll want to know: \ur{ہانکنا}}
\textbf{\ur{ہانکنا}} — \emph{hāṅknā} — \textquotedblleft{}to drive (animals or a cart)\textquotedblright{}.

\begin{quote}
\ur{ایک} \ur{کسان} \ur{شاہراہ} \ur{کے} \ur{کنارے} \ur{بکریاں} \ur{ہانک} \ur{رہا} \ur{تھا۔}
\end{quote}

\begin{quote}
\emph{ek kisān shāhrāh ke kināre bakriyāṅ hāṅk rahā thā.}
\end{quote}

\begin{quote}
A farmer was driving goats along the edge of the highway.
\end{quote}

Past, with \ur{نے} (he or she …): \ur{اس} \ur{نے} \ur{گھوڑا} \ur{ہانکا}, \emph{us ne ghoṛā hāṅkā}.

\begin{cousinweb}
\textbf{\ur{ہانکنا}} (\emph{hāṅknā}) is inherited: CDIAL lists the Hindi cousin \emph{hā̃knā}, \textquotedblleft{}to cry out, to drive off\textquotedblright{}, under Sanskrit \emph{hakkayati}, \textquotedblleft{}calls to\textquotedblright{} (CDIAL 13939): animals are driven along with shouts.
\end{cousinweb}

\begin{grammarlens}[title={What you've built}]
One more word for getting around.
\end{grammarlens}

\subsection*{Your turn}
\begin{itemize}
\raggedright
  \item \emph{Say it:} \emph{hāṅknā}
  \item \emph{Say it:} \emph{hāṅknā}, once more
  \item \emph{Say it:} \emph{ek kisān shāhrāh ke kināre bakriyāṅ hāṅk rahā thā.} — A farmer was driving goats along the edge of the highway.
  \item \emph{Recall:} say the Urdu for a checkpoint (on a road), then the Urdu for a convoy, then say \emph{hāṅknā} again
\end{itemize}

\begin{tcolorbox}[breakable,colback=teal!4,colframe=teal!35!black,title={Before you move on}]
What does \ur{ہانکنا} mean? (\textquotedblleft{}To drive (animals or a cart)\textquotedblright{}.) Say it once more.
\end{tcolorbox}
